\chapter{Conclusion and Future Work}
\label{sec:conclusion}

Identifying cancer driver genes (CDGs) accurately remains a critical challenge in cancer research. These genes play a crucial role in cancer development and progression. Despite extensive work in the field, predicting CDGs remains complex due to the heterogeneity of underlying biological mechanisms. In recent years, the use of graph-based learning methods, particularly Graph Neural Networks (GNNs), has shown great potential to improve the accuracy of CDG prediction. \citet{schulte2021integration}, for example, produced an approach capable of integrating pan-cancer data into a GCN-based model to identify new drivers of cancer. Our study expanded on this approach by evaluating and comparing different methodological strategies to predict CDGs using three different GNN algorithms. We focused on adding useful information to describe genes and dealing with class imbalance. Centralities added relevant information to node features vector, increasing performance in graph-based learning methods as well as for traditional ML algorithms. We also observed that graph-based learning algorithms generally outperformed other approaches in the prediction task.

Among the applied GNNs, GCN demonstrated the most promising results, showing consistent and robust performance in all evaluated scenarios. We carried out an extensive search for hyperparameter optimization to optimize prediction results in each of the six PPI networks selected for the study.
Moreover, we successfully implemented two ensemble learning strategies, a majority voting system, and a system using the average predicted clas probabilities among models networks. The overall performance of these models in predicting test set genes showed a clear impact of the two ensemble approaches in reducing the number of false positives. We compared these results with a comprehesinve PPI network defined in this work, which combines the nodes and edges of all six selected PPI networks and removes duplicate gene links. The ensemble models outperformed the other approaches used for prediction, surpassing both the individual PPI networks and the unified PPI network. Our best approach, represented by the Ensemble-Average GCN-based model, reached an AUC-PR score of 0.67 in an independent test set of drivers and passenger genes. Although this result is promising, our model still failed to improve over the performance of previously proposed approaches \cite{schulte2021integration}.

Further research in this area could lead to a better understanding of the molecular mechanisms underlying cancer and pave the way for more precise models that, in turn, will help in the development of more effective diagnostic and therapeutic strategies. Areas for growth and development include new approaches for handling class imbalance in graph-based learning, such as GraphSMOTE \cite{zhao2021graphsmote}, and multi-task learning on GNNs to be able to simultaneously train models capable of predicting cancer-specific and pan-cancer driver genes. Understanding driver mutation profiles of different cancer types is crucial, as driver mutations can differ between cancer types, and individual tumors have their own mutation profiles. Finally, improving the strategies to create labeled data for positive and negative examples of CDGs can have a great impact on the overall performance of prediction model. 


